\documentclass[11pt,a4paper]{article}
\usepackage[margin=27mm]{geometry}
\usepackage{amsmath,amssymb,amsthm,hyperref}
\newtheorem{theorem}{Theorem}\newtheorem{proposition}[theorem]{Proposition}
\newtheorem{lemma}[theorem]{Lemma}
\newcommand{\E}{\mathbb E}\newcommand{\Pp}{\mathbb P}
\title{Four summands: moment paths and active collisions}
\author{}\date{28 October 2026}
\begin{document}\maketitle

A linear moment theorem does not compress an iid four-fold convolution.
The counterexample below occurs inside the bounded moment and threshold
box used by the current four-summand reduction. A contact inequality does
supply a different restriction: apart from at most two points, every atom
of a finite maximizing law must occur in an active four-sum collision.
This is a structural reduction with explicit remaining cases, not a proof
of the four-summand Berry--Esseen bound.

\section{An exact obstruction to moment-extreme compression}
Let $Z_\theta$ have support $\{-2,-1,0,1,2\}$ and probability vector
\begin{equation}\label{eq:family}
 \left(\frac1{20}+\theta,\ \frac15-2\theta,\ \frac12,
              \frac15+2\theta,\ \frac1{20}-\theta\right),
 \qquad -\frac1{20}\le\theta\le\frac1{20}.
\end{equation}
All these laws have
\[
 \E Z_\theta=0,\qquad \E Z_\theta^2=\frac45,
 \qquad \E|Z_\theta|^3=\frac65.
\]
Consequently $X_\theta=(\sqrt5/2)Z_\theta$ is standardized and has
$\rho=3\sqrt5/4<46/25$. Its support is contained in
$[-\sqrt5,\sqrt5]\subset(-4.032,4.032)$, and
$\beta=\E[X_\theta|X_\theta|]=-5\theta$, so $|\beta|\le1/4$.
The fixed standardized raw sum threshold corresponding to
$\sum_{j=1}^4Z_{\theta,j}\le1$ is $H=\sqrt5/2$.

\begin{proposition}\label{prop:obstruction}
On the set of probability laws supported on these five raw points with
these three moments, the four-summand closed CDF at threshold $1$ attains
its maximum strictly inside the moment segment. In particular neither
extreme law maximizes it. The same assertion holds for the standardized
fixed-target functional
\[
 2\{\Pp(X_{\theta,1}+\cdots+X_{\theta,4}\le H)-\Phi(H/2)\}
                          -c\rho
\]
for every fixed $c$, because its other terms are constant on the segment.
\end{proposition}
\begin{proof}
The matrix with rows $1,z,z^2,|z|^3$ on the five support points has rank
four. Its nullspace is spanned by $(1,-2,0,2,-1)$; nonnegativity gives
exactly the interval in \eqref{eq:family}. Hence the two endpoints are
the only extreme points of this moment polytope.
Direct convolution, with all threshold atoms included, gives
\begin{align}\label{eq:cdf}
 F(\theta)&:=\Pp(Z_{\theta,1}+\cdots+Z_{\theta,4}\le1)\notag\\
 &=\frac{25719}{32000}-\frac{103}{2000}\theta
       -\frac{483}{200}\theta^2-\frac{123}{5}\theta^3+11\theta^4.
\end{align}
In particular
\[
 F(0)=\frac{25719}{32000},\quad
 F(-1/20)=\frac{4017}{5000},\quad
 F(1/20)=\frac{7921}{10000}.
\]
The two strict differences are $51/160000$ and $1859/160000$.
Continuity gives a maximizer, and neither endpoint can be one.
More explicitly,
\[
 F'(\theta)=\frac{(20\theta+1)(4400\theta^2-7600\theta-103)}{2000}.
\]
The second factor decreases on the entire interval and changes sign
between $-1/40$ and $-1/100$. Thus the unique maximum is at
\[
 \theta_* =\frac{190-\sqrt{37233}}{220}\in(-1/40,-1/100).
\]
The standardization and fixed-target assertions follow from the common
moments and common threshold.
\end{proof}

This is not a counterexample to the Berry--Esseen conjecture. It refutes
an intermediate extremal principle. It even has a positive closed
normal discrepancy: $F(0)>4/5$, whereas
$\Phi(\sqrt5/4)\le1/2+\sqrt5/10<73/100$.
The example does not preserve the signed third moment:
$\E X_\theta^3=-(15\sqrt5/2)\theta$. A reduction imposing that additional
constraint would require its own argument. Nor does the example assert
that its interior maximizer is a global fixed-target maximizer over all
probability laws.

\section{The correct finite moment-path calculation}
For finite signed measures define
\[
 \mathcal A_H(\mu_1,\mu_2,\mu_3,\mu_4)
  =\int\mathbf1_{\{x_1+x_2+x_3+x_4\le H\}}
                         \,d\mu_1\,d\mu_2\,d\mu_3\,d\mu_4.
\]
This is symmetric and multilinear. The notation
$\mathcal A_H(P^{4-j},\nu^j)$ records how many arguments are of each type.

\begin{lemma}[A complete one-dimensional replacement]\label{lem:path}
Let $P$ be a probability measure with finite absolute third moment and
let $\nu$ be a finite signed measure with
$\int(1+|x|^3)\,d|\nu|(x)<\infty$. Assume $P+t\nu$ is a nonnegative
probability measure for $t\in[a,b]$. Suppose
\[
 \int(1,x,x^2,|x|^3)\,d\nu(x)=0.
\]
The mean, variance and absolute third moment remain fixed, and
\begin{equation}\label{eq:expansion}
 F_{P+t\nu}^{*4}(H)=
 \sum_{j=0}^4\binom4j t^j\mathcal A_H(P^{4-j},\nu^j).
\end{equation}
Its maximum is attained at an endpoint or an interior zero of the cubic
 derivative. Thus an exact derivative-root check, or a whole-interval
quartic majorant, is sufficient. Checking only the endpoints is not.
\end{lemma}
\begin{proof}
The moments are linear. Expand the four product factors, grouping equal
terms by symmetry. The event stays fixed, including its boundary; no
continuity of its indicator is used. The remaining assertion is the
ordinary extreme-value theorem and Fermat's theorem for a polynomial.
\end{proof}

A useful necessary test for a local maximum is also immediate. If both
signs of $t$ are feasible near zero, then
\begin{gather*}
 \mathcal A_H(P^3,\nu)=0,\qquad
 \mathcal A_H(P^2,\nu^2)\le0.
\end{gather*}
If the second quantity vanishes, then also
$\mathcal A_H(P,\nu^3)=0$ and $\mathcal A_H(\nu^4)\le0$.
For $P=\sum_{j=1}^m p_j\delta_{x_j}$ with $p_j>0$, write
$\nu=\sum_jv_j\delta_{x_j}$ and
\[
 K_{ij}=P^{*2}(( -\infty,H-x_i-x_j]).
\]
The quadratic condition is $v^\mathsf TKv\le0$ for every vector in
the nullspace of the four moment rows. In other words, the compression
of $K$ to that nullspace must be negative semidefinite. This is a necessary
same-law condition; it does not characterize an optimizer or replace the
whole event probability. It offers an inexpensive rejection test for
candidate atomic supports before any dense certificate construction.

\section{At most two atoms can avoid the active collision}
We now use a contact premise, which is distinct from a moment-polytope
extreme-point claim.

\begin{lemma}[A cubic contact potential has at most two minima]
\label{lem:minima}
For $c>0$ and real $a,b,d$, the function
$q(x)=c|x|^3+ax^2+bx+d$ has at most two distinct local minima.
If there are two, one is negative and one positive.
\end{lemma}
\begin{proof}
If $a\ge0$, the function is strictly convex, so has at most one local
minimum. If $a<0$, set $R=-a/(3c)>0$. On $(-R,R)$ its second
derivative $6c|x|+2a$ is negative. On each exterior interval it is
strictly positive, so the derivative has at most one zero there.
A stationary point at $R$ or $-R$ has a nonzero third derivative and
is not a local minimum. These observations give at most one local
minimum on each exterior interval and none elsewhere.
\end{proof}

\begin{theorem}[Contact forces active support collisions]\label{thm:collision}
Let $P$ have finite support $S$, with every listed atom of positive mass.
Fix $H\in\mathbb R$, $c>0$, and assume that for some real $a,b,d$
\begin{equation}\label{eq:contact}
 q(x):=c|x|^3+ax^2+bx+d
 \ \ge\ 8P^{*3}(( -\infty,H-x]),
 \qquad x\in\mathbb R,
\end{equation}
with equality at every $x\in S$. Then at most two points of $S$ fail
to belong to an active four-sum collision. Equivalently,
\[
 \#\{x\in S:H-x\notin S+S+S\}\le2.
\]
If there are two such exceptional points, they have opposite signs.
If exactly one unordered four-term multiset from $S$ sums to $H$, and
it contains $r$ distinct points, then $\#S\le r+2\le6$.
\end{theorem}
\begin{proof}
If $H-x\notin S+S+S$, the finite-support CDF
$P^{*3}(( -\infty,H-y])$ is constant for $y$ in a neighborhood of
$x$. Equality at $x$ and \eqref{eq:contact} therefore make $x$ a
local minimum of the same function $q$. Apply Lemma~\ref{lem:minima}.
Every point that does occur in an active collision belongs to its
four-term multiset. If there is only one such multiset, it contains
all but at most two support points, which proves the last assertion.
\end{proof}

For the fixed-target objective
$2\{P^{*4}(( -\infty,H])-\Phi(H/2)\}-c\rho$, the contact formula
in \texttt{sections/02-fixed-target.tex} has exactly this form.
For a standardized maximizing pair, put
$\varphi=\phi(H/2)$, $F=P^{*4}(( -\infty,H])$ and
$\beta=\E[X|X|]$. Its potential is
\[
 q(x)=c|x|^3-\frac{H\varphi+3c\rho}{2}x^2
              -(4\varphi+3c\beta)x
              +8F+\frac{H\varphi+c\rho}{2}.
\]
The fixed-target contact theorem supplies \eqref{eq:contact}; it is an
explicit premise of the new support conclusion, not proved by a finite
arithmetic check. No two- or three-summand estimate is used here.
For a finite law, a maximizing threshold must itself be a four-sum atom:
otherwise lowering it a little preserves the CDF and strictly improves
the normal subtraction.

In particular, if distinct unordered four-term multisets of $S$ always
have distinct sums, a finite fixed-target maximizing law has at most six
atoms. Conversely, a finite maximizing law with at least seven atoms must
have two distinct active occupancy vectors $k,l\in\mathbb N^m$, with
\[
 \sum_jk_j=\sum_jl_j=4,\qquad
 \sum_j(k_j-l_j)x_j=0,\qquad
 0<\sum_j|k_j-l_j|\le8.
\]
Thus larger atomic candidates must lie on explicit low-integer-coefficient
resonance hyperplanes. The two possible exceptional atoms can be restricted
further to the local minima of the displayed cubic potential.

\paragraph{What remains open.}
This is not a universal six-atom reduction. Attainment of the unrestricted
problem gives compact support, not finite support. A finite maximizing law
may have several simultaneous active multisets; these resonant strata
cannot be discarded as a set of parameter measure zero. Perturbing a
threshold collision changes a closed event discontinuously. Nor does a
finite atomic approximation of a maximizing law inherit its exact contact
inequality. A complete short $n=4$ proof must address these remaining cases,
or use an original-law bound such as the retained physical ROOT4 bridge.

\paragraph{Scope of the bounded checks.}
The independent rational script verifies the moment segment, its complete
quartic by both convolution and direct product enumeration, its endpoint
and derivative identities, moment-null tangent expansions, and actual
collision hypergraphs. It rejects endpoint-only maximization and the
omission of the contact premise. These are exact controls for the lemmas,
not a new acceptance of G2, G3, ROOT4 or the finite base through $41$.
\end{document}
